\scanpage{70}
Notons dans ce contexte le résultat plus profond de Igusa

\uline{Proposition}. $X$ comme ci-dessus, \deleted{\unsure{soit}} posons
\[
E(\overset{\inserted{\unsure{k|1}}}{X})=\sum_{i=0}^{4}(-1)^i b_i(X)
\]
avec
\[
\left\{\begin{aligned}
b_0(X)&=b_4(X)=1\\
b_1(X)&=b_3(X)=\dim\underline{\mathrm{Pic}}_X\\
E(X)&=c_2(X)=\text{self-intersection de la diagonale dans } X\times X.
\end{aligned}\right.
\]
ce qui détermine donc $b_2(X)$ de façon unique, \inserted{$\boxed{b_2=E-2+2p}$ où $p=\dim\underline{\mathrm{Pic}}_X$} Ceci dit, posons
\[
\rho=\operatorname{rang}(\underline{\text{NérSév}}_{X/k}),
\]
alors on a l'inégalité
\[
\boxed{\rho\leq b_2}
\]
% archive sheet number: 69

